\begin{figure}[htbp]
\centering
\begin{tikzpicture}[x=.58cm,y=.58cm]
 \foreach \dx/\sx/\sy in {0/0/0,5.8/1/0,11.6/0/1,17.4/1/1}{
  \begin{scope}[shift={(\dx,0)}]
   \bcuadricula{4}
   \draw[btinta,line width=.9pt] (1.5,3.5)--(3.5,1.5);
   \ifnum\numexpr\sx+\sy\relax=2\relax
    \draw[btinta,dash pattern=on 2pt off 1.5pt,line width=.9pt] (1.5,1.5)--(3.5,3.5);
   \fi
   \foreach \fr/\fc/\libre in {1/2/1,3/4/1,1/4/\sx,3/2/\sy}{
    \begin{scope}[shift={(\fc-.5,4.5-\fr)}]
     \fill[bpapel] (0,0) circle[radius=.26];
     \ifnum\libre=1\relax
      \draw[btinta,line width=1pt] (0,0) circle[radius=.22];
     \else
      \draw[btinta,line width=1.05pt] (-.17,-.17)--(.17,.17) (-.17,.17)--(.17,-.17);
     \fi
    \end{scope}
    \ifnum\libre=1\relax \bmarca{4}{\fr}{\fc}\fi
   }
   \node[anchor=north,font=\small] at (2,-.18) {$(X,Y)=(\sx,\sy)$};
  \end{scope}
 }
\end{tikzpicture}
\caption{Los cuatro estados del cuadrado de frontera se dibujan en una ventana de
\(4\times4\) de un tablero mayor, con coordenadas locales \(r=1\), \(q=2\), \(h=2\).
El círculo indica una casilla disponible y la cruz una vetada. La línea continua une
\(e=(1,2)\) y \(f=(3,4)\). La discontinua une \((1,4)\) con \((3,2)\) cuando ambas
están disponibles. Las demás casillas se muestran sólo como referencia. En el estado \(11\),
tras restar los efectos de cada par, los coeficientes de interacción en el logaritmo del conteo
y en el del predictor son \(-1\) y \(+2\), en unidades del producto de los cuatro pesos al orden
principal. Su diferencia \(-1-(+2)=-3\) da el signo de \(-S\) en~\eqref{eq:exceso-cuadrado}.}
\label{fig:frontera}
\end{figure}
